\begin{center}\begin{tikzpicture}[font=\small,>=stealth,line join=round]
\draw[gray!40,rounded corners=3pt] (0,0) rectangle (14.2,8.3);\node at (7.1,7.700000000000001) {例1：内偶先合并，外层加三不改变同高};
\begin{scope}[shift={(4.1,2.161111111111111)},x=1.25cm,y=0.5111111111111111cm]\draw[->,gray] (-2.4,0)--(2.4,0) node[right]{$x$};\draw[->,gray] (0,-1)--(0,8) node[above]{$y$};\draw[blue,thick,solid,domain=-2.1:2.1,samples=180] plot(\x,{(\x)^2+3});\fill[orange!85!black] (-2,7) circle(2.2pt);\fill[orange!85!black] (2,7) circle(2.2pt);\end{scope}\node[anchor=west,align=left,text width=5.0cm] at (8.4,4.0) {$F=x^2+3$\\[10pt]$g=x^2$，$f(u)=u+3$\\[10pt]$F(2)=F(-2)=7$\\[10pt]外层本身非奇非偶};
\node[text width=12.6cm,align=center] at (7.1,.65) {向上平移不破坏左右同高；这不是要求外层也具有奇偶性。};
\end{tikzpicture}\end{center}
